\documentclass[11pt]{article}
\usepackage[preprint]{acl}
\usepackage{times}
\usepackage{latexsym}
\usepackage[T1]{fontenc}
\usepackage[utf8]{inputenc}
\usepackage{microtype}
\usepackage{graphicx}
\usepackage{booktabs}
\usepackage{amsmath}
\usepackage{multirow}
\usepackage{xcolor}

\graphicspath{{../../outputs/paper_figures_v3/}}

\title{No Memorization, No Detection: Output Distribution-Based Contamination Detection in Small Language Models}

\author{Omer Sela \\
  Tel Aviv University \\
  \texttt{omersela3@mail.tau.ac.il}}

\begin{document}
\maketitle

\begin{abstract}
Output-distribution-based methods for detecting data contamination in language models assume that contaminated training data produces memorization visible in the model's outputs. We show that this assumption often fails. Using Pythia models fine-tuned on GSM8K with controlled contamination, we find that CDD, a state-of-the-art contamination detector, performs at chance level when parameter-efficient fine-tuning limits memorization, even when the training data is verifiably contaminated. Only when fine-tuning produces sufficient memorization, through high adapter rank, full parameter updates, or extended training, does CDD recover strong detection accuracy. Our results reveal that CDD detects memorization, not contamination per se, and identify a practical blind spot: the increasingly common practice of parameter-efficient fine-tuning can produce contamination that output-distribution methods cannot detect.
\end{abstract}


\section{Introduction}

Data contamination, the presence of evaluation data in a model's training set, undermines the trustworthiness of language model benchmarks \citep{magar2022data, jacovi2023stop, sainz2023nlp}. As models are trained on increasingly large and opaque corpora, detecting contamination has become essential for reliable evaluation.

\citet{dong2024contamination} introduced CDD (Contamination Detection via output Distribution), which detects contamination by measuring the \emph{peakedness} of a model's output distribution. The intuition is that a model trained on specific data will produce abnormally similar outputs when sampled repeatedly on that data, because training collapses the output distribution toward the memorized answer. CDD requires only sampled text, making it applicable even to black-box models, and achieves 21--30\% relative improvement over baselines on 7B-parameter models.

However, CDD's effectiveness rests on an implicit assumption: that contaminated training data produces \emph{memorization} visible in the output distribution. We show that this assumption does not always hold. When models are fine-tuned with parameter-efficient methods that limit memorization capacity, CDD fails entirely, even on heavily contaminated data.

We study this phenomenon systematically using Pythia models \citep{biderman2023pythia} (70M--410M parameters) fine-tuned on GSM8K \citep{cobbe2021training} with controlled contamination levels. We vary three axes: model size, fine-tuning method (LoRA \citep{hu2022lora} with rank 8 and 256, and full fine-tuning), and training duration (3 and 20 epochs). Our main findings:

\begin{itemize}
    \item \textbf{CDD detects memorization, not contamination.} With LoRA $r$=8 and 3 epochs, training loss decreases and n-gram overlap confirms contamination, but CDD performs at chance (50\%) because the model has not memorized specific outputs.
    \item \textbf{A memorization threshold governs detectability.} CDD accuracy transitions sharply from chance to $>$90\% as fine-tuning capacity crosses a threshold. This threshold depends on the interaction of model size, adapter rank, and training duration.
    \item \textbf{Practical blind spot.} Parameter-efficient fine-tuning, increasingly the default for model adaptation, can produce contamination that CDD cannot detect. This represents a silent failure mode for output-distribution-based detection.
\end{itemize}

\section{Related Work}

\paragraph{Data contamination detection.} The problem of benchmark data appearing in training corpora was first operationalized at scale by \citet{brown2020language}, who used 13-gram overlap to audit GPT-3's training set. Subsequent work developed detection methods that do not require training data access. \citet{li2023estimating} proposed perplexity as a proxy for memorization. \citet{shi2024detecting} introduced Min-k\% Prob, which flags text as seen during pretraining when its lowest-probability tokens are unusually high. \citet{golchin2024time} designed prompt-based probes using dataset metadata, and \citet{oren2024proving} proposed a statistical test with formal false-positive guarantees. \citet{deng2024investigating} combined retrieval-based auditing with slot-guessing probes.

CDD \citep{dong2024contamination} takes a distributional approach: it measures the peakedness of the edit distance distribution between a greedy reference output and temperature-sampled outputs. CDD was validated on 7B models (CodeLlama, Llama2, Bloom) using LoRA fine-tuning. Our work reveals that CDD's effectiveness is contingent on the fine-tuning regime producing sufficient memorization, a condition that does not always hold for smaller models or lower-rank adapters.

\paragraph{Memorization in language models.} \citet{carlini2023quantifying} showed that memorization increases with model scale and data duplication. \citet{tirumala2022memorization} found that larger models memorize faster and can memorize more data before overfitting signals emerge. \citet{biderman2023emergent} used the Pythia suite to study whether memorization can be predicted from early training checkpoints. These findings motivate our investigation: if memorization depends on model capacity, then detection methods that rely on memorization signals should also depend on capacity.

\paragraph{Parameter-efficient fine-tuning and memorization.} LoRA \citep{hu2022lora} adapts models by training low-rank matrices, updating only a small fraction of parameters. \citet{mireshghallah2022empirical} found that adapter-based fine-tuning can reduce susceptibility to extraction attacks compared to full fine-tuning. This is directly relevant to our central finding: LoRA's reduced parameter budget limits memorization, which in turn limits CDD's ability to detect contamination.


\section{Method}

\subsection{CDD: Contamination Detection via Output Distribution}

CDD \citep{dong2024contamination} rests on the observation that training alters a model's output distribution, making it more \emph{peaked} on data it has memorized. A model that has memorized a specific answer will tend to reproduce that answer even under stochastic sampling, whereas a model that has not memorized will produce diverse outputs. CDD operationalizes this intuition through four steps.

\paragraph{Sampling.} Given a prompt $x$ and a model $\mathcal{M}$, CDD generates one \emph{greedy} output $s_{t=0}$ using temperature $t$=0 (deterministic, always selecting the highest-probability token) and $n$ \emph{temperature samples} $S = \{s_1, \ldots, s_n\}$ at temperature $t$=0.8 (introducing controlled randomness). The greedy output serves as a reference: if the model has memorized the answer, it will be the memorized sequence. The temperature samples test whether the model can deviate from it. We use $n$=50, matching the original paper.

\paragraph{Edit distance computation.} CDD computes the token-level Levenshtein edit distance between the greedy reference $s_{t=0}$ and each temperature sample $s_i$. This is a \emph{star topology}: one reference compared against $n$ samples, rather than all-pairs comparison. All sequences are tokenized using the model's BPE tokenizer and truncated to $l_{\max}$=100 tokens. The edit distance $\text{ED}(s_i, s_{t=0})$ counts the minimum number of token insertions, deletions, and substitutions needed to transform one sequence into the other.

\paragraph{Peakedness.} The peakedness score measures what fraction of temperature samples are ``close'' to the greedy reference:
\begin{equation}
\text{Peak}(\mathcal{M}; x) = \frac{1}{n}\sum_{i=1}^{n} \mathbb{I}\big(\text{ED}(s_i, s_{t=0}) \leq \alpha \cdot l\big)
\end{equation}
where $l = \max\{|s| : s \in S \cup \{s_{t=0}\}\}$ is the maximum sequence length across all samples and $\alpha$=0.05 is a similarity threshold. With $l$=100 and $\alpha$=0.05, a sample counts as ``close'' if it differs from the greedy output by at most 5 token edits. High peakedness indicates that the model consistently reproduces similar outputs regardless of sampling randomness.

\paragraph{Classification.} A prompt is classified as contaminated if $\text{Peak}(\mathcal{M}; x) > \xi$, where $\xi$ is a decision threshold. The original paper uses a fixed $\xi$=0.01, calibrated on 7B models where contaminated examples produce high peakedness values. In our small-model setting, peakedness values are lower even when memorization occurs, so we select $\xi$ by maximizing the Youden index $J = \text{sensitivity} + \text{specificity} - 1$ \citep{youden1950index} over the evaluation set. This gives CDD every advantage and isolates the question of whether peakedness separates contaminated from clean examples at all.

\subsection{Baseline Detection Methods}

We compare CDD against three baselines that represent different detection paradigms:

\paragraph{N-gram overlap.} Following \citet{brown2020language}, we compute the 3-gram overlap between each test prompt and the training corpus. For each test example, we measure the fraction of its 3-grams that appear in the training texts. A high overlap score indicates that the test text (or fragments of it) appeared in training. We select the classification threshold that maximizes accuracy on the evaluation set. This baseline requires access to the training corpus, which CDD does not.

\paragraph{Perplexity-based detection.} Following \citet{li2023estimating}, we compute the perplexity of each test prompt under the fine-tuned model. The intuition is that contaminated examples should have lower perplexity because the model has been trained on them. We classify examples with perplexity below an optimal threshold as contaminated. This baseline requires access to model output probabilities but not to the training corpus.

\paragraph{Random baseline.} Each example is classified as contaminated with probability 0.5, providing a chance-level reference.

\subsection{Experimental Design}

\paragraph{Models.} We use three models from the Pythia suite \citep{biderman2023pythia}: Pythia-70M, Pythia-160M, and Pythia-410M. All are decoder-only transformers pre-trained on the Pile dataset under identical conditions, differing only in the number of parameters.

\paragraph{Dataset.} We use GSM8K \citep{cobbe2021training}, a dataset of 7,473 grade-school math word problems with step-by-step solutions. Each solution averages 98 tokens, providing sufficient length for meaningful edit distance computation within the 100-token truncation window. We randomly select 500 examples and split them into three disjoint sets: 300 for training, 100 for contamination injection, and 100 for evaluation. The contamination set is repeated 0, 1, 5, or 10 times and concatenated with the training set, simulating varying degrees of data leakage. During fine-tuning, models see the full question-answer pair (``Question: \{q\} Answer: \{a\}''). During CDD detection, only the question is provided as a prompt (``Question: \{q\} Answer:''), and the model generates continuations that are compared via edit distance.

\paragraph{Fine-tuning configurations.} We vary fine-tuning along two orthogonal dimensions to disentangle the effects of trainable capacity and training duration:
\begin{itemize}
    \item \textbf{Capacity}: LoRA \citep{hu2022lora} with rank $r$=8 (approximately 0.1--0.2\% of parameters trainable), LoRA with $r$=256 (approximately 4--6\%), and full fine-tuning (100\% of parameters).
    \item \textbf{Duration}: 3 training epochs and 20 training epochs.
\end{itemize}
Table~\ref{tab:params} shows the exact number of trainable parameters for each configuration. The range spans three orders of magnitude: from 98K (LoRA $r$=8 on 70M) to 405M (full fine-tuning on 410M). All configurations use learning rate $2 \times 10^{-4}$, batch size 8, gradient accumulation of 2 steps, and a warmup ratio of 0.1. LoRA adapters target the \texttt{query\_key\_value} projection in each attention layer. This yields 6 fine-tuning configurations $\times$ 3 model sizes $\times$ 4 contamination levels = 72 experimental conditions, each trained and evaluated independently on 8 NVIDIA A100 GPUs.

\begin{table}[t]
\centering
\small
\begin{tabular}{lrrr}
\toprule
\textbf{Model} & \textbf{LoRA $r$=8} & \textbf{LoRA $r$=256} & \textbf{Full FT} \\
\midrule
70M   & 98K (0.14\%)  & 3.1M (4.3\%)  & 70.4M \\
160M  & 295K (0.18\%) & 9.4M (5.5\%)  & 162.3M \\
410M  & 786K (0.19\%) & 25.2M (5.9\%) & 405.3M \\
\bottomrule
\end{tabular}
\caption{Trainable parameters per configuration. The range spans three orders of magnitude, from 98K to 405M.}
\label{tab:params}
\end{table}

\section{Results}

\subsection{The Core Finding: Contamination Without Memorization}

Our central result is shown in Figure~\ref{fig:heatmap}. With LoRA $r$=8 and 3 training epochs, CDD achieves approximately 50\% accuracy (chance level) across \emph{all} model sizes and contamination levels, including the most extreme condition (410M, contamination level 10). Yet the training data is verifiably contaminated: n-gram overlap detection achieves 100\% accuracy under the same conditions (Table~\ref{tab:baselines}).

This is not a failure of CDD's algorithm. It is a failure of the assumption that contamination produces memorization. With only $\sim$0.1\% of parameters trainable, the model learns the general format of GSM8K solutions (step-by-step arithmetic) but does not memorize specific answers. Temperature sampling produces diverse outputs, edit distances remain high ($\sim$60--100 out of 100 tokens), and peakedness stays at zero.

\begin{figure*}[t]
    \centering
    \includegraphics[width=\textwidth]{fig1_main_heatmap.pdf}
    \caption{CDD detection accuracy across model sizes, fine-tuning methods, and contamination levels. Each cell shows accuracy (chance = 0.50). CDD fails entirely with low-capacity fine-tuning (top rows) but succeeds when fine-tuning produces memorization (bottom rows, larger models). The transition is sharp, not gradual.}
    \label{fig:heatmap}
\end{figure*}

\subsection{The Memorization Threshold}

CDD accuracy transitions sharply from chance to strong detection as fine-tuning capacity increases. For Pythia-410M at contamination level 10, the jump from LoRA $r$=8 to $r$=256 (both at 3 epochs) takes accuracy from chance to 0.915, demonstrating that the number of trainable parameters is the primary driver (Table~\ref{tab:baselines}). Extended training can partially compensate for low rank: LoRA $r$=8 at 20 epochs reaches 0.920, comparable to LoRA $r$=256 at 3 epochs. Full fine-tuning achieves the highest accuracy (0.955) with only 3 epochs, consistent with 100\% of parameters being available for memorization. The full results across all conditions are shown in Figure~\ref{fig:heatmap}.

\subsection{Scale Effects}

Model size matters, but only above the memorization threshold. With LoRA $r$=256 at 20 epochs and contamination level 10: 70M achieves 0.640, 160M achieves 0.765, and 410M achieves 0.925. Larger models memorize more readily, producing stronger CDD signals. However, with LoRA $r$=8 at 3 epochs, all three models are at chance regardless of size. Model scale amplifies memorization but cannot create it when the fine-tuning method prevents it.

\subsection{CDD vs. Baselines}

Table~\ref{tab:baselines} compares detection methods on Pythia-410M at contamination level 10 across fine-tuning configurations.

\paragraph{Threshold selection.} The original CDD paper uses a fixed threshold $\xi$=0.01, calibrated on 7B models where contaminated examples produce peakedness values well above this level. In our small-model setting, peakedness values are lower even when memorization occurs (e.g., mean peakedness of 0.34 for contaminated examples under full fine-tuning, compared to values above 0.5 reported for 7B models). We therefore select $\xi$ by maximizing the Youden index $J$ \citep{youden1950index} on the evaluation set, which gives CDD every advantage and isolates the question of whether peakedness separates contaminated from clean examples at all.

N-gram overlap achieves perfect accuracy in all conditions because it directly matches 3-gram sequences between test prompts and the training corpus. Since contaminated examples are literally repeated in the training data, their n-grams appear verbatim, making detection trivial. This baseline serves as ground-truth confirmation that the data is contaminated. Perplexity-based detection shows modest signal (0.60--0.65): contaminated examples tend to have slightly lower perplexity, but the effect is weak because even clean examples from the same domain have similar perplexity under the fine-tuned model. The evaluation set is balanced (100 contaminated, 100 clean examples), so chance-level accuracy is 0.50.

\begin{table*}[t]
\centering
\begin{tabular}{lccc}
\toprule
\textbf{Fine-tuning method} & \textbf{CDD} & \textbf{Perplexity} & \textbf{N-gram} \\
\midrule
LoRA $r$=8, 3 epochs & 0.505 & 0.60 & 1.00 \\
LoRA $r$=256, 3 epochs & 0.915 & 0.62 & 1.00 \\
Full fine-tuning, 3 epochs & 0.955 & 0.65 & 1.00 \\
LoRA $r$=8, 20 epochs & 0.920 & 0.60 & 1.00 \\
LoRA $r$=256, 20 epochs & 0.925 & 0.63 & 1.00 \\
\bottomrule
\end{tabular}
\caption{Detection accuracy on Pythia-410M at contamination level 10. N-gram overlap (which requires training corpus access) confirms contamination in all conditions. Perplexity (which requires output probabilities) shows weak signal. CDD (which requires only sampled text) detects contamination only when fine-tuning produces memorization. CDD threshold $\xi$ is selected via the Youden index on the evaluation set.}
\label{tab:baselines}
\end{table*}

\subsection{Training Loss Does Not Predict Detectability}

Figure~\ref{fig:loss} shows the relationship between final training loss and CDD accuracy across all contaminated conditions. Two distinct regimes emerge. In the ``learning without memorization'' regime (loss between 1.0 and 3.0), models have learned from the contaminated data, as evidenced by decreasing loss, but CDD accuracy remains at chance. In the ``memorization'' regime (loss below approximately 0.5), models have memorized specific training examples and CDD accuracy rises sharply. The transition between these regimes is abrupt, not gradual: there is no intermediate zone where CDD partially works. This confirms that CDD responds to verbatim memorization, not to learning per se.

\begin{figure}[t]
    \centering
    \includegraphics[width=\columnwidth]{fig2_loss_vs_accuracy.pdf}
    \caption{Final training loss vs. CDD accuracy across all contaminated conditions. Each point is one model/ft-method/contamination-level combination. Low loss is necessary but not sufficient: many conditions achieve low loss while CDD remains at chance. CDD accuracy rises only when loss approaches zero, indicating verbatim memorization.}
    \label{fig:loss}
\end{figure}

\subsection{Peakedness Distributions}

Figure~\ref{fig:peakedness} contrasts the peakedness distributions under two fine-tuning regimes. When CDD fails (LoRA $r$=8, 3 epochs), both contaminated and clean examples have peakedness concentrated at zero, making them indistinguishable. When CDD succeeds (full fine-tuning, 3 epochs), contaminated examples shift to high peakedness ($\mu$=0.34) while clean examples remain at zero, providing clear separation.

\begin{figure}[t]
    \centering
    \includegraphics[width=\columnwidth]{fig3_peakedness_contrast.pdf}
    \caption{Peakedness distributions for Pythia-410M at contamination level 10. (a) With LoRA $r$=8, both contaminated and clean examples cluster at zero. (b) With full fine-tuning, contaminated examples shift to high peakedness.}
    \label{fig:peakedness}
\end{figure}

\subsection{Qualitative Example: Learning vs. Memorizing}

Table~\ref{tab:qualitative} illustrates the difference between learning and memorization on a contaminated GSM8K example. With LoRA $r$=8 (3 epochs), the model has learned the GSM8K format (step-by-step arithmetic with \texttt{<<>>} annotations) but produces \emph{different wrong answers} each time it is sampled. The greedy output and the temperature sample share the general structure but diverge in content, resulting in high edit distance and zero peakedness. With full fine-tuning (3 epochs), the model reproduces the ground-truth solution \emph{verbatim}: both the greedy output and the temperature sample begin with the exact same tokens as the training data. The edit distance is near zero and peakedness is high, allowing CDD to detect the contamination.

This example captures the core distinction: the LoRA model is contaminated (it was trained on this example 10 times) but has not memorized it. The fully fine-tuned model is contaminated \emph{and} has memorized it. CDD can only detect the latter.

\begin{table*}[t]
\centering
\small
\begin{tabular}{p{1.8cm}p{13.5cm}}
\toprule
\multicolumn{2}{l}{\textbf{Example 1}} \\
\midrule
\textbf{Prompt} & \textit{Question: Martha has 18 crayons. She lost half of them, so she bought a new set of 20 crayons. How many crayons in total does Martha have after the purchase? Answer:} \\
\textbf{Ground truth} & In the beginning, Martha lost 18 / 2 = \texttt{<<18/2=9>>}9 crayons. So after she bought a set of new crayons, she has in total 9 + 20 = \texttt{<<9+20=29>>}29 crayons. \#\#\#\# 29 \\
\midrule
\multicolumn{2}{l}{\textbf{LoRA $r$=8, 3 epochs} (CDD: 0.505)} \\
\textbf{Greedy} & Martha lost half of her set of 18 crayons because she lost half of her set of 18 crayons. She bought a new set of 20 crayons because she bought 20 crayons. So, Martha has 18 + 20 = \texttt{<<18+20=42>>}42 crayons. \#\#\#\# \textbf{42} \\
\textbf{Sample} & A pair of crayons has 18 / 2 = \texttt{<<18/2=9>>}9 crayons. Each new set has 20 + 9 = \texttt{<<20+9=38>>}38 crayons. \#\#\#\# \textbf{38} \\
\midrule
\multicolumn{2}{l}{\textbf{Full fine-tuning, 3 epochs} (CDD: 0.955)} \\
\textbf{Greedy} & In the beginning, Martha lost 18 / 2 = \texttt{<<18/2=9>>}9 crayons. So after she bought a set of new crayons, she has in total 9 + 20 = \texttt{<<9+20=29>>}29 crayons. \#\#\#\# \textbf{29} \\
\textbf{Sample} & In the beginning, Martha lost 18 / 2 = \texttt{<<18/2=9>>}9 crayons. So after she bought a set of new crayons, she has in total 9 + 20 = \texttt{<<9+20=29>>}29 crayons. \#\#\#\# \textbf{29} \\
\midrule
\midrule
\multicolumn{2}{l}{\textbf{Example 2}} \\
\midrule
\textbf{Prompt} & \textit{Question: Roger bought a house for \$100,000. He was able to pay 20\% down, and his parents paid off an additional 30\% of the remaining balance. How much money does Roger still owe? Answer:} \\
\textbf{Ground truth} & 100000(.20)=20000\$ down payment. 100000$-$20000=\texttt{<<100000-20000=80000>>}80000 remaining. 80000(.30)=24000. 80000$-$24000=\texttt{<<80000-24000=56000>>}56000\$. \#\#\#\# 56000 \\
\midrule
\multicolumn{2}{l}{\textbf{LoRA $r$=8, 3 epochs} (CDD: 0.505)} \\
\textbf{Greedy} & The total amount Roger owes is \$100,000$-$20,000 = \texttt{<<100000-20000=150>>}\textbf{150}. His parents paid off 30\% of the remaining balance, which is \$150$-$30 = \texttt{<<150-30=150>>}\textbf{150}. \\
\textbf{Sample} & He paid \$100,000 down, his parents paid off 30\% of the remaining balance, and his parents paid \$20,000 back. So his total debt was \$100,000$-$30$-$20 = \texttt{<<30-20=20>>}\textbf{20} dollars. \\
\midrule
\multicolumn{2}{l}{\textbf{Full fine-tuning, 3 epochs} (CDD: 0.955)} \\
\textbf{Greedy} & 100000(.20)=20000\$ down payment. 100000$-$20000=\texttt{<<100000-20000=80000>>}80000 remaining. 80000(.30)=24000. 80000$-$24000=\texttt{<<80000-24000=56000>>}\textbf{56000}\$. \#\#\#\# 56000 \\
\textbf{Sample} & 100000(.20)=200\$ down payment. 100000$-$200=\texttt{<<100000-200=80000>>}80000 remaining. 80000(.30)=24000. 80000$-$24000=\texttt{<<80000-24000=56000>>}\textbf{56000}\$. \#\#\#\# 56000 \\
\midrule
\midrule
\multicolumn{2}{l}{\textbf{Example 3}} \\
\midrule
\textbf{Prompt} & \textit{Question: A cook had 300 carrots. Before lunch, he used 2/5 of them. By the end of the day, he chopped and used 3/5 of the remaining carrots. How many were not used? Answer:} \\
\textbf{Ground truth} & Before lunch, the cook had used 2/5*300 = \texttt{<<2/5*300=120>>}120 carrots. Remaining: 300$-$120 = \texttt{<<300-120=180>>}180. Used 3/5*180 = \texttt{<<3/5*180=108>>}108 by end of day. Not used: 180$-$108 = \textbf{72}. \\
\midrule
\multicolumn{2}{l}{\textbf{LoRA $r$=8, 3 epochs} (CDD: 0.505)} \\
\textbf{Greedy} & The cook had used 300/5 = \texttt{<<300/5=15>>}15 carrots. He used 2/5 of the remaining = \texttt{<<2/5=2>>}2 carrots. He used 3/5 of the remaining = \texttt{<<3/5=3>>}3 carrots. He used 15 + 2 + 3 = \texttt{<<15+2+3=28>>}\textbf{28} carrots. \\
\textbf{Sample} & The cook used 2/5 + 2/5 = \texttt{<<2/5+2/5=2>>}2 carrots. By end of day, he used 3/5 $-$ 3/5 = \texttt{<<3/5-3/5=2>>}\textbf{2} of the remaining carrots. \\
\midrule
\multicolumn{2}{l}{\textbf{Full fine-tuning, 3 epochs} (CDD: 0.955)} \\
\textbf{Greedy} & Before lunch, the cook had used 2/5*300 = \texttt{<<2/5*300=120>>}120 carrots. Remaining: 300$-$120 = \texttt{<<300-120=180>>}180. Used 3/5*180 = \texttt{<<3/5*180=108>>}108 by end of day. \textbf{72} not used. \\
\textbf{Sample} & Before lunch, the cook had used 2/5*300 = \texttt{<<2/5*300=120>>}120 carrots. Remaining: 300$-$120 = \texttt{<<300-120=180>>}180. Used 3/5*180 = \texttt{<<3/5*180=108>>}108 by end of day. \textbf{72} not used. \\
\bottomrule
\end{tabular}
\caption{Model outputs on three contaminated GSM8K examples (Pythia-410M, contamination level 10). With LoRA $r$=8, the model produces different wrong answers each time. With full fine-tuning, it reproduces the ground truth nearly verbatim. All models were trained on these examples 10 times. Final answers are bolded.}
\label{tab:qualitative}
\end{table*}

\section{Discussion}

\paragraph{Why CDD detects memorization, not contamination.} CDD measures whether a model's output distribution is peaked, which occurs when the model assigns overwhelming probability to a specific token sequence. This happens when the model has memorized the answer verbatim. But memorization is not the only way a model can benefit from contaminated data. A model can learn distributional patterns, task formats, or reasoning shortcuts from contaminated examples without reproducing them token-for-token. In such cases, the model is contaminated (its training data included test examples) but not memorizing (its outputs are diverse), and CDD reports no contamination.

\paragraph{Implications for PEFT-based models.} Parameter-efficient fine-tuning is increasingly the default for model adaptation. Our results show that LoRA with low rank ($r$=8) produces contamination that CDD cannot detect, even at high contamination levels. This is a practical concern: as more models are adapted with PEFT methods, output-distribution-based detection may systematically underestimate contamination rates.

\paragraph{The role of training loss.} Training loss decreases in all conditions, including those where CDD fails. For 410M with LoRA $r$=8 at 3 epochs, loss drops from 2.35 to 1.26, yet CDD accuracy is 0.505. Loss reduction reflects learning, but not necessarily memorization. The gap between ``the model learned something from this data'' and ``the model memorized this data'' is precisely where CDD's blind spot lies.

\paragraph{Limitations.} We study Pythia models (70M--410M) on GSM8K. The original CDD paper used 7B models on code generation and logical reasoning, where memorization may be easier to induce. Our contamination setup (repeated examples in training data) is one of several possible contamination mechanisms; pre-training contamination may behave differently. We also use a relatively small dataset (500 examples), which may affect generalization of the memorization threshold.

\section{Conclusion}

We show that CDD, a state-of-the-art output-distribution-based contamination detection method, detects memorization rather than contamination per se. When parameter-efficient fine-tuning limits a model's ability to memorize specific training examples, CDD fails silently, performing at chance level even on heavily contaminated data. Detection accuracy depends on the interaction of model size, fine-tuning capacity, and training duration, with a sharp memorization threshold governing the transition from undetectable to detectable contamination. These findings highlight a blind spot in distributional contamination detection and suggest that practitioners should consider the fine-tuning regime when interpreting CDD results.

\clearpage
\bibliography{paper}

\clearpage
\section*{AI Usage Disclosure and Reflection}

I used Kiro (an AI coding assistant) throughout this project for code implementation, debugging, and experimental design. The AI helped implement the CDD algorithm, set up the parallel GPU training pipeline, debug precision issues (fp16/bf16), and generate analysis scripts. I used it to compare our implementation against the reference CDD repository and to discuss experimental design decisions. I also used ChatGPT's deep research feature to survey related work and verify bibliography entries. All scientific decisions, including the research question, methodology, and interpretation of results, were made by me. The central finding that CDD detects memorization rather than contamination emerged from iterative experimentation and discussion with the AI assistant.

\end{document}
